% =========================================================
% Chapter 8: Linear Momentum and Collisions
% POSN 1 Physics Preparation Notes
% Language: Thai
% Compiler: XeLaTeX
% =========================================================

\chapter{โมเมนตัมเชิงเส้นและการชน}

\begin{center}
    {\Large \textbf{Linear Momentum and Collisions}}\\[4pt]
    {\normalsize เอกสารเตรียมสอบคัดเลือก สอวน. ฟิสิกส์ ค่าย 1}\\[6pt]
    {\normalsize จัดทำโดย \textbf{Tames Thanakrit Damduan}}\\[2pt]
    {\footnotesize \textcopyright\ 2026 Tames Thanakrit Damduan. All rights reserved.}\\[8pt]
\end{center}

\section*{เป้าหมายของบทเรียน}

หลังจากเรียนบทนี้ นักเรียนควรทำสิ่งต่อไปนี้ได้

\begin{enumerate}
    \item เข้าใจความหมายของโมเมนตัมเชิงเส้นในฐานะปริมาณเวกเตอร์
    \item เขียนกฎข้อที่สองของนิวตันในรูปโมเมนตัมได้
    \item เข้าใจแรงดลและดล รวมถึงความสัมพันธ์ระหว่างดลกับการเปลี่ยนโมเมนตัม
    \item ใช้กฎอนุรักษ์โมเมนตัมในหนึ่งมิติและสองมิติได้
    \item แยกชนยืดหยุ่น ชนไม่ยืดหยุ่น และชนติดกันได้
    \item ใช้พลังงานจลน์ร่วมกับโมเมนตัมในโจทย์ชนยืดหยุ่นได้
    \item ใช้แนวคิดศูนย์กลางมวลและความเร็วศูนย์กลางมวลได้
    \item แก้โจทย์ระเบิด การดีดตัว การชนกับสปริง และโจทย์ชนระดับ สอวน. ได้
\end{enumerate}

\begin{tcolorbox}[title=แนวคิดสำคัญของบทนี้]
โมเมนตัมคือภาษาของการเคลื่อนที่เมื่อมีการชน การระเบิด หรือแรงกระทำในช่วงเวลาสั้น ๆ จุดสำคัญคือถ้าแรงภายนอกลัพธ์ต่อระบบเป็นศูนย์ โมเมนตัมรวมของระบบจะคงตัว แม้พลังงานจลน์อาจไม่คงตัวก็ตาม
\end{tcolorbox}

\newpage{}

% =========================================================
\section{โมเมนตัมเชิงเส้น}

\subsection{นิยามของโมเมนตัม}

วัตถุที่มีมวลมากและเคลื่อนที่เร็ว ย่อมหยุดหรือเปลี่ยนทิศได้ยากกว่าวัตถุที่เบาและเคลื่อนที่ช้า ปริมาณที่บอก ``ความยากในการเปลี่ยนสภาพการเคลื่อนที่'' นี้คือโมเมนตัม

โมเมนตัมเชิงเส้นนิยามเป็น

\[
\vec{p}=m\vec{v}
\]

โดยที่

\[
m=\text{มวลของวัตถุ}
\]

และ

\[
\vec{v}=\text{ความเร็วของวัตถุ}
\]

\begin{tcolorbox}[title=นิยาม]
โมเมนตัมเชิงเส้นของวัตถุคือ

\[
\vec{p}=m\vec{v}
\]

โมเมนตัมเป็นเวกเตอร์ และมีทิศเดียวกับความเร็ว
\end{tcolorbox}

\subsection{หน่วยของโมเมนตัม}

จากนิยาม

\[
\vec{p}=m\vec{v}
\]

หน่วยของโมเมนตัมคือ

\[
\si{kg.m/s}
\]

และเพราะ

\[
1\,\si{N}=1\,\si{kg.m/s^2}
\]

จึงมี

\[
1\,\si{N.s}=1\,\si{kg.m/s}
\]

ดังนั้นหน่วยของโมเมนตัมอาจเขียนได้ทั้ง

\[
\si{kg.m/s}
\]

หรือ

\[
\si{N.s}
\]

\subsection{โมเมนตัมเป็นเวกเตอร์}

ถ้าเลือกแกนขวาเป็นบวก วัตถุมวล $m$ เคลื่อนที่ไปทางขวาด้วยอัตราเร็ว $v$ จะมี

\[
p=+mv
\]

ถ้าเคลื่อนที่ไปทางซ้ายด้วยอัตราเร็ว $v$ จะมี

\[
p=-mv
\]

\begin{tcolorbox}[title=ข้อควรจำ]
อย่าคิดโมเมนตัมเป็นแค่ $mv$ แบบไม่มีทิศทาง ในโจทย์ชน เครื่องหมายของโมเมนตัมสำคัญมาก
\end{tcolorbox}

\begin{examplebox}{ตัวอย่าง: โมเมนตัมในหนึ่งมิติ}
วัตถุมวล $2\,\si{kg}$ เคลื่อนที่ไปทางขวาด้วยอัตราเร็ว $5\,\si{m/s}$ จงหาโมเมนตัม

เลือกขวาเป็นบวก

\[
p=mv
\]

\[
p=(2)(5)
\]

\[
p=10\,\si{kg.m/s}
\]

ดังนั้น

\[
\boxed{p=+10\,\si{kg.m/s}}
\]
\end{examplebox}

\begin{exercisebox}{แบบฝึกหัด 8.1}
\begin{enumerate}
    \item วัตถุมวล $3\,\si{kg}$ เคลื่อนที่ด้วยความเร็ว $4\,\si{m/s}$ ไปทางขวา จงหาโมเมนตัม
    \item วัตถุมวล $0.5\,\si{kg}$ เคลื่อนที่ด้วยความเร็ว $-20\,\si{m/s}$ จงหาโมเมนตัม
    \item วัตถุสองก้อนมีมวลเท่ากัน แต่ก้อนแรกมีอัตราเร็วเป็นสองเท่าของก้อนที่สอง จงเปรียบเทียบโมเมนตัม
    \item วัตถุสองก้อนมีอัตราเร็วเท่ากัน แต่ก้อนแรกมีมวลเป็นสามเท่าของก้อนที่สอง จงเปรียบเทียบโมเมนตัม
\end{enumerate}
\end{exercisebox}

% =========================================================
\section{กฎข้อที่สองของนิวตันในรูปโมเมนตัม}

\subsection{เริ่มจากกรณีมวลคงตัว}

กฎข้อที่สองของนิวตันเขียนได้ว่า

\[
\sum \vec{F}=m\vec{a}
\]

ถ้ามวลคงตัว และ

\[
\vec{a}=\frac{d\vec{v}}{dt}
\]

จะได้

\[
\sum \vec{F}=m\frac{d\vec{v}}{dt}
\]

เพราะ $m$ คงตัว จึงเขียนเป็น

\[
\sum \vec{F}=\frac{d(m\vec{v})}{dt}
\]

แต่

\[
m\vec{v}=\vec{p}
\]

ดังนั้น

\[
\sum \vec{F}=\frac{d\vec{p}}{dt}
\]

\begin{tcolorbox}[title=กฎข้อที่สองในรูปโมเมนตัม]
\[
\sum \vec{F}=\frac{d\vec{p}}{dt}
\]

แรงลัพธ์คืออัตราการเปลี่ยนโมเมนตัมตามเวลา
\end{tcolorbox}

\subsection{ความหมาย}

ถ้าแรงลัพธ์มาก โมเมนตัมจะเปลี่ยนเร็ว

ถ้าแรงลัพธ์เป็นศูนย์ โมเมนตัมไม่เปลี่ยน

\[
\sum \vec{F}=0
\quad \Rightarrow \quad
\frac{d\vec{p}}{dt}=0
\]

ดังนั้น

\[
\vec{p}=\text{ค่าคงตัว}
\]

นี่คือจุดเริ่มของกฎอนุรักษ์โมเมนตัม

% =========================================================
\section{ดลและแรงดล}

\subsection{นิยามของดล}

ในสถานการณ์ที่แรงกระทำในช่วงเวลาสั้นมาก เช่น ลูกบอลชนพื้น ค้อนตอกตะปู หรือวัตถุชนกัน เรามักไม่รู้รายละเอียดของแรงทุกขณะ แต่รู้ว่ามันกระทำในช่วงเวลา $\Delta t$

ดลของแรงนิยามเป็น

\[
\vec{J}=\int_{t_i}^{t_f}\vec{F}\,dt
\]

ถ้าแรงเฉลี่ยเป็น $\vec{F}_{\text{avg}}$ จะได้

\[
\vec{J}=\vec{F}_{\text{avg}}\Delta t
\]

\begin{tcolorbox}[title=ดล]
\[
\vec{J}=\int \vec{F}\,dt
\]

ถ้าใช้แรงเฉลี่ย

\[
\vec{J}=\vec{F}_{\text{avg}}\Delta t
\]
\end{tcolorbox}

\subsection{ทฤษฎีบทดลและโมเมนตัม}

เริ่มจาก

\[
\sum \vec{F}=\frac{d\vec{p}}{dt}
\]

จัดรูป

\[
\sum \vec{F}\,dt=d\vec{p}
\]

อินทิเกรตจากเวลาเริ่มต้นถึงเวลาสุดท้าย

\[
\int_{t_i}^{t_f}\sum \vec{F}\,dt
=
\int_{\vec{p}_i}^{\vec{p}_f}d\vec{p}
\]

ด้านซ้ายคือดลรวม

\[
\vec{J}_{\text{net}}
\]

ด้านขวาคือการเปลี่ยนโมเมนตัม

\[
\vec{p}_f-\vec{p}_i
\]

ดังนั้น

\[
\vec{J}_{\text{net}}=\Delta \vec{p}
\]

\begin{tcolorbox}[title=ทฤษฎีบทดลและโมเมนตัม]
\[
\vec{J}_{\text{net}}=\Delta \vec{p}
\]

หรือ

\[
\vec{F}_{\text{avg}}\Delta t=m\vec{v}_f-m\vec{v}_i
\]
\end{tcolorbox}

\subsection{แรงเฉลี่ยระหว่างการชน}

จาก

\[
\vec{F}_{\text{avg}}\Delta t=\Delta \vec{p}
\]

จะได้

\[
\vec{F}_{\text{avg}}=\frac{\Delta \vec{p}}{\Delta t}
\]

ถ้าเวลาในการชนสั้น แรงเฉลี่ยจะมาก

\begin{tcolorbox}[title=ข้อควรจำ]
ถ้าเปลี่ยนโมเมนตัมเท่าเดิม แต่เพิ่มเวลาที่ใช้ในการเปลี่ยนโมเมนตัม แรงเฉลี่ยจะลดลง หลักนี้อธิบายได้ว่าทำไมถุงลมนิรภัย เบาะรอง และการงอเข่าตอนลงพื้นจึงช่วยลดแรงกระแทก
\end{tcolorbox}

\begin{examplebox}{ตัวอย่าง: แรงเฉลี่ยจากดล}
ลูกบอลมวล $0.20\,\si{kg}$ เคลื่อนที่ไปทางขวาด้วยความเร็ว $10\,\si{m/s}$ แล้วชนกำแพงและสะท้อนกลับด้วยความเร็ว $6\,\si{m/s}$ ไปทางซ้าย ถ้าเวลาชนคือ $0.04\,\si{s}$ จงหาแรงเฉลี่ยที่กำแพงกระทำต่อลูกบอล

เลือกขวาเป็นบวก

\[
v_i=+10\,\si{m/s}
\]

\[
v_f=-6\,\si{m/s}
\]

การเปลี่ยนโมเมนตัมคือ

\[
\Delta p=m(v_f-v_i)
\]

\[
\Delta p=(0.20)(-6-10)
\]

\[
\Delta p=-3.2\,\si{kg.m/s}
\]

ดังนั้นแรงเฉลี่ยคือ

\[
F_{\text{avg}}=\frac{\Delta p}{\Delta t}
\]

\[
F_{\text{avg}}=\frac{-3.2}{0.04}
\]

\[
F_{\text{avg}}=-80\,\si{N}
\]

เครื่องหมายลบแปลว่าแรงเฉลี่ยชี้ไปทางซ้าย

\[
\boxed{80\,\si{N} \text{ ไปทางซ้าย}}
\]
\end{examplebox}

\begin{exercisebox}{แบบฝึกหัด 8.2}
\begin{enumerate}
    \item วัตถุมวล $1\,\si{kg}$ เปลี่ยนความเร็วจาก $2\,\si{m/s}$ เป็น $8\,\si{m/s}$ ในทิศเดียวกัน จงหาดล
    \item ลูกบอลมวล $0.5\,\si{kg}$ ตกลงพื้นด้วยความเร็ว $10\,\si{m/s}$ แล้วหยุดในเวลา $0.02\,\si{s}$ จงหาขนาดแรงเฉลี่ยจากพื้น โดยไม่คิดน้ำหนักระหว่างชน
    \item ลูกบอลสองลูกมีมวลเท่ากันและตกจากความสูงเท่ากัน ลูกหนึ่งกระดอนสูงกว่าอีกลูก ถ้าเวลาสัมผัสพื้นเท่ากัน ลูกใดได้รับแรงเฉลี่ยจากพื้นมากกว่า
    \item ถ้าดลคงเดิม แต่เวลาชนเพิ่มเป็นสองเท่า แรงเฉลี่ยเปลี่ยนอย่างไร
\end{enumerate}
\end{exercisebox}

% =========================================================
\section{การอนุรักษ์โมเมนตัม}

\subsection{ระบบอนุภาค}

ให้ระบบประกอบด้วยอนุภาคหลายตัว โมเมนตัมรวมของระบบคือผลรวมแบบเวกเตอร์ของโมเมนตัมทุกตัว

\[
\vec{P}=\vec{p}_1+\vec{p}_2+\vec{p}_3+\cdots
\]

หรือ

\[
\vec{P}=\sum_i m_i\vec{v}_i
\]

\begin{tcolorbox}[title=โมเมนตัมรวมของระบบ]
\[
\vec{P}=\sum_i \vec{p}_i
\]
\end{tcolorbox}

\subsection{เงื่อนไขของการอนุรักษ์โมเมนตัม}

ถ้าแรงภายนอกลัพธ์ที่กระทำต่อระบบเป็นศูนย์

\[
\sum \vec{F}_{\text{ext}}=0
\]

จะได้ว่าโมเมนตัมรวมคงตัว

\[
\vec{P}_i=\vec{P}_f
\]

\begin{tcolorbox}[title=กฎอนุรักษ์โมเมนตัม]
ถ้าแรงภายนอกลัพธ์ต่อระบบเป็นศูนย์ หรือมีค่าน้อยมากในช่วงเวลาที่พิจารณา

\[
\vec{P}_i=\vec{P}_f
\]

หรือ

\[
\sum \vec{p}_i=\sum \vec{p}_f
\]
\end{tcolorbox}

\subsection{แรงภายในไม่เปลี่ยนโมเมนตัมรวม}

ในการชน วัตถุสองก้อนออกแรงต่อกัน แต่แรงที่วัตถุหนึ่งกระทำต่ออีกวัตถุหนึ่งเป็นคู่แรงกิริยาและปฏิกิริยาตามกฎข้อที่สามของนิวตัน

แรงภายในเหล่านี้อาจเปลี่ยนโมเมนตัมของวัตถุแต่ละก้อน แต่ผลรวมของโมเมนตัมทั้งระบบไม่เปลี่ยน ถ้าไม่มีแรงภายนอกลัพธ์

\begin{tcolorbox}[title=ข้อควรจำสำหรับโจทย์ชน]
ระหว่างการชน แรงภายในมักมีค่ามากกว่าแรงภายนอกมาก ดังนั้นในช่วงเวลาสั้น ๆ ของการชน เรามักถือว่าโมเมนตัมรวมของระบบคงตัว
\end{tcolorbox}

% =========================================================
\section{การชนในหนึ่งมิติ}

\subsection{ชนติดกัน}

ชนติดกันคือการชนที่หลังชนวัตถุเคลื่อนที่ไปด้วยกันด้วยความเร็วเดียวกัน

ให้มวล $m_1$ มีความเร็วต้น $u_1$ และมวล $m_2$ มีความเร็วต้น $u_2$

หลังชนติดกัน ความเร็วร่วมคือ $v$

ใช้อนุรักษ์โมเมนตัม

\[
m_1u_1+m_2u_2=(m_1+m_2)v
\]

ดังนั้น

\[
v=\frac{m_1u_1+m_2u_2}{m_1+m_2}
\]

\begin{tcolorbox}[title=ชนติดกัน]
\[
v=\frac{m_1u_1+m_2u_2}{m_1+m_2}
\]

ชนติดกันเป็นชนไม่ยืดหยุ่นอย่างสมบูรณ์ พลังงานจลน์ไม่คงตัว
\end{tcolorbox}

\begin{examplebox}{ตัวอย่าง: ชนติดกัน}
มวล $2\,\si{kg}$ เคลื่อนที่ไปทางขวาด้วย $6\,\si{m/s}$ เข้าชนมวล $4\,\si{kg}$ ที่อยู่นิ่ง แล้วติดกันไป จงหาความเร็วหลังชน

เลือกขวาเป็นบวก

\[
m_1u_1+m_2u_2=(m_1+m_2)v
\]

แทนค่า

\[
(2)(6)+(4)(0)=(2+4)v
\]

\[
12=6v
\]

\[
v=2\,\si{m/s}
\]

ดังนั้น

\[
\boxed{v=2\,\si{m/s} \text{ ไปทางขวา}}
\]
\end{examplebox}

\subsection{ชนไม่ยืดหยุ่น}

ชนไม่ยืดหยุ่นคือการชนที่โมเมนตัมรวมคงตัว แต่พลังงานจลน์ไม่คงตัว

\[
\vec{P}_i=\vec{P}_f
\]

แต่โดยทั่วไป

\[
K_i\neq K_f
\]

พลังงานจลน์บางส่วนเปลี่ยนเป็นความร้อน เสียง การเสียรูป หรือพลังงานภายใน

\subsection{ชนยืดหยุ่น}

ชนยืดหยุ่นคือการชนที่ทั้งโมเมนตัมรวมและพลังงานจลน์รวมคงตัว

\[
\vec{P}_i=\vec{P}_f
\]

และ

\[
K_i=K_f
\]

\begin{tcolorbox}[title=ชนยืดหยุ่น]
ชนยืดหยุ่นต้องใช้ทั้งสองสมการ

\[
m_1u_1+m_2u_2=m_1v_1+m_2v_2
\]

\[
\frac{1}{2}m_1u_1^2+\frac{1}{2}m_2u_2^2
=
\frac{1}{2}m_1v_1^2+\frac{1}{2}m_2v_2^2
\]
\end{tcolorbox}

\subsection{สูตรชนยืดหยุ่นในหนึ่งมิติ}

สำหรับการชนยืดหยุ่นในหนึ่งมิติ

\[
v_1=
\frac{m_1-m_2}{m_1+m_2}u_1
+
\frac{2m_2}{m_1+m_2}u_2
\]

\[
v_2=
\frac{2m_1}{m_1+m_2}u_1
+
\frac{m_2-m_1}{m_1+m_2}u_2
\]

\begin{tcolorbox}[title=ข้อควรระวัง]
สูตรนี้ใช้เฉพาะการชนยืดหยุ่นในหนึ่งมิติเท่านั้น ถ้าโจทย์เป็นสองมิติ หรือลูกกลมชนเฉียง ต้องแตกความเร็วเป็นแนวตามเส้นศูนย์กลางและแนวสัมผัสก่อน
\end{tcolorbox}

\subsection{กรณีมวลที่สองอยู่นิ่ง}

ถ้า $u_2=0$ จะได้

\[
v_1=\frac{m_1-m_2}{m_1+m_2}u_1
\]

\[
v_2=\frac{2m_1}{m_1+m_2}u_1
\]

ถ้า $m_1=m_2$

\[
v_1=0
\]

\[
v_2=u_1
\]

คือมวลแรกหยุด และมวลที่สองรับความเร็วไปทั้งหมด

\begin{examplebox}{ตัวอย่าง: ชนยืดหยุ่นหนึ่งมิติ}
นิวตรอนมวล $m$ เคลื่อนที่ด้วยความเร็ว $u$ ชนตรง ๆ แบบยืดหยุ่นกับดิวเทอรอนมวล $2m$ ที่อยู่นิ่ง จงหาอัตราส่วนพลังงานจลน์ของนิวตรอนหลังชนต่อก่อนชน

ใช้สูตรชนยืดหยุ่นเมื่อมวลที่สองอยู่นิ่ง

\[
v_1=\frac{m_1-m_2}{m_1+m_2}u_1
\]

แทน

\[
m_1=m,\quad m_2=2m,\quad u_1=u
\]

\[
v_1=\frac{m-2m}{m+2m}u
\]

\[
v_1=-\frac{1}{3}u
\]

พลังงานจลน์หลังชนของนิวตรอนคือ

\[
K_f=\frac{1}{2}m\left(\frac{u}{3}\right)^2
\]

พลังงานจลน์ก่อนชนคือ

\[
K_i=\frac{1}{2}mu^2
\]

ดังนั้น

\[
\frac{K_f}{K_i}
=
\frac{\frac{1}{2}m\frac{u^2}{9}}{\frac{1}{2}mu^2}
\]

\[
\boxed{\frac{K_f}{K_i}=\frac{1}{9}}
\]
\end{examplebox}

\begin{exercisebox}{แบบฝึกหัด 8.3}
\begin{enumerate}
    \item มวล $m$ เคลื่อนที่ด้วยความเร็ว $u$ ชนติดกับมวล $2m$ ที่อยู่นิ่ง จงหาความเร็วหลังชน
    \item มวล $m$ เคลื่อนที่ด้วยความเร็ว $u$ ชนยืดหยุ่นในหนึ่งมิติกับมวล $m$ ที่อยู่นิ่ง จงหาความเร็วหลังชนของทั้งสองก้อน
    \item มวล $m$ เคลื่อนที่ด้วยความเร็ว $u$ ชนยืดหยุ่นกับมวล $3m$ ที่อยู่นิ่ง จงหาความเร็วหลังชนของมวลแรก
    \item ในการชนติดกัน พลังงานจลน์คงตัวหรือไม่ จงอธิบาย
\end{enumerate}
\end{exercisebox}

% =========================================================
\section{ศูนย์กลางมวล}

\subsection{ตำแหน่งศูนย์กลางมวล}

สำหรับระบบอนุภาคหลายตัว ตำแหน่งศูนย์กลางมวลคือค่าเฉลี่ยถ่วงน้ำหนักด้วยมวล

ในหนึ่งมิติ

\[
x_{\text{cm}}=\frac{m_1x_1+m_2x_2+\cdots}{m_1+m_2+\cdots}
\]

ในรูปทั่วไป

\[
\vec{r}_{\text{cm}}=
\frac{\sum_i m_i\vec{r}_i}{\sum_i m_i}
\]

\begin{tcolorbox}[title=ศูนย์กลางมวล]
\[
\vec{r}_{\text{cm}}=\frac{\sum_i m_i\vec{r}_i}{M_{\text{tot}}}
\]
\end{tcolorbox}

\subsection{ความเร็วศูนย์กลางมวล}

หาอนุพันธ์ของตำแหน่งศูนย์กลางมวล

\[
\vec{v}_{\text{cm}}=\frac{\sum_i m_i\vec{v}_i}{M_{\text{tot}}}
\]

แต่

\[
\sum_i m_i\vec{v}_i=\vec{P}_{\text{tot}}
\]

ดังนั้น

\[
\vec{P}_{\text{tot}}=M_{\text{tot}}\vec{v}_{\text{cm}}
\]

\begin{tcolorbox}[title=โมเมนตัมรวมกับศูนย์กลางมวล]
\[
\vec{P}_{\text{tot}}=M_{\text{tot}}\vec{v}_{\text{cm}}
\]
\end{tcolorbox}

\subsection{ใช้ศูนย์กลางมวลกับสปริงและการชน}

ในระบบที่ไม่มีแรงภายนอกในแนวที่พิจารณา ความเร็วศูนย์กลางมวลคงตัว

จุดนี้สำคัญมากในโจทย์ที่มวลหนึ่งชนสปริงที่ติดกับอีกมวลหนึ่ง เพราะขณะที่สปริงถูกอัดมากที่สุด วัตถุทั้งสองเข้าใกล้กันไม่ได้อีกแล้ว จึงมีความเร็วเท่ากันชั่วขณะ

ความเร็วนั้นคือความเร็วศูนย์กลางมวล

\[
v_{\text{common}}=v_{\text{cm}}
\]

\begin{examplebox}{ตัวอย่าง: มวลชนสปริงที่ติดกับอีกมวลหนึ่ง}
มวล $m$ เคลื่อนที่ด้วยความเร็ว $u$ เข้าชนสปริงที่ติดกับมวล $M$ ซึ่งอยู่นิ่งบนพื้นลื่น จงหาความเร็วของมวล $m$ ขณะที่สปริงถูกอัดมากที่สุด

เมื่อสปริงถูกอัดมากที่สุด มวล $m$ และ $M$ มีความเร็วเท่ากันชั่วขณะ ให้เป็น $V$

ใช้อนุรักษ์โมเมนตัม

\[
mu+M(0)=(m+M)V
\]

ดังนั้น

\[
\boxed{V=\frac{mu}{m+M}}
\]
\end{examplebox}

% =========================================================
\section{การชนและการระเบิดในสองมิติ}

\subsection{อนุรักษ์โมเมนตัมแยกแกน}

โมเมนตัมเป็นเวกเตอร์ ถ้าไม่มีแรงภายนอกลัพธ์ ต้องอนุรักษ์โมเมนตัมทั้งแกน $x$ และแกน $y$

\[
P_{xi}=P_{xf}
\]

\[
P_{yi}=P_{yf}
\]

\begin{tcolorbox}[title=การอนุรักษ์โมเมนตัมในสองมิติ]
\[
\sum p_{xi}=\sum p_{xf}
\]

\[
\sum p_{yi}=\sum p_{yf}
\]
\end{tcolorbox}

\subsection{การระเบิด}

การระเบิดคือสถานการณ์ที่วัตถุหนึ่งก้อนแตกออกเป็นหลายชิ้นด้วยแรงภายในระบบ ถ้าแรงภายนอกในช่วงเวลาระเบิดมีผลน้อยมาก โมเมนตัมรวมก่อนและหลังระเบิดยังคงเท่ากัน

\begin{examplebox}{ตัวอย่าง: ระเบิดเป็นสองชิ้น}
มวล $m$ เคลื่อนที่ด้วยความเร็ว $u$ ตามแกน $+x$ แล้วระเบิดออกเป็นสองชิ้น ชิ้นแรกมีมวล $\alpha m$ และเคลื่อนที่ด้วยความเร็ว $v$ ตามแกน $+y$ จงหาทิศของชิ้นที่สอง

โมเมนตัมก่อนระเบิดคือ

\[
\vec{p}_i=mu\hat{i}
\]

หลังระเบิด ชิ้นแรกมีโมเมนตัม

\[
\vec{p}_1=\alpha mv\hat{j}
\]

ชิ้นที่สองมีมวล

\[
(1-\alpha)m
\]

ให้ความเร็วของชิ้นที่สองเป็น

\[
\vec{v}_2=v_{2x}\hat{i}+v_{2y}\hat{j}
\]

ใช้อนุรักษ์โมเมนตัม

\[
mu\hat{i}
=
\alpha mv\hat{j}
+
(1-\alpha)m(v_{2x}\hat{i}+v_{2y}\hat{j})
\]

แยกแกน $x$

\[
mu=(1-\alpha)mv_{2x}
\]

\[
v_{2x}=\frac{u}{1-\alpha}
\]

แยกแกน $y$

\[
0=\alpha mv+(1-\alpha)mv_{2y}
\]

\[
v_{2y}=-\frac{\alpha v}{1-\alpha}
\]

ดังนั้นชิ้นที่สองเคลื่อนที่ลงใต้แกน $x$

ถ้าให้ $\theta$ เป็นมุมที่ชิ้นที่สองทำกับแกน $+x$ ลงด้านล่าง

\[
\tan\theta=\frac{|v_{2y}|}{v_{2x}}
\]

\[
\boxed{\tan\theta=\frac{\alpha v}{u}}
\]
\end{examplebox}

\subsection{ชนเฉียงของทรงกลมผิวลื่น}

ถ้าทรงกลมผิวลื่นชนกัน แรงดลระหว่างชนจะอยู่ตามแนวเส้นที่เชื่อมศูนย์กลางของทรงกลมทั้งสองในขณะชน

ดังนั้นให้แตกความเร็วเป็นสองแนว

\[
\text{แนวเส้นศูนย์กลาง}
\]

และ

\[
\text{แนวสัมผัส}
\]

สำหรับทรงกลมผิวลื่น แรงดลไม่มีองค์ประกอบในแนวสัมผัส ดังนั้นความเร็วแนวสัมผัสของวัตถุแต่ละก้อนจะไม่เปลี่ยน

ในกรณีมวลเท่ากัน และก้อนที่สองเริ่มนิ่ง ถ้าชนยืดหยุ่น องค์ประกอบความเร็วตามแนวเส้นศูนย์กลางจะถ่ายไปยังก้อนที่สอง ส่วนก้อนแรกเหลือเฉพาะองค์ประกอบแนวสัมผัส

\begin{tcolorbox}[title=กลยุทธ์สำหรับชนเฉียง]
\begin{enumerate}
    \item วาดแนวเส้นศูนย์กลางขณะชน
    \item แตกความเร็วเป็นแนวเส้นศูนย์กลางและแนวสัมผัส
    \item ใช้สมการชนหนึ่งมิติเฉพาะตามแนวเส้นศูนย์กลาง
    \item แนวสัมผัสไม่เปลี่ยนถ้าผิวลื่นและไม่มีแรงเสียดทาน
\end{enumerate}
\end{tcolorbox}

% =========================================================
\section{พลังงานในการชน}

\subsection{พลังงานจลน์ที่สูญเสียในการชนติดกัน}

สำหรับการชนติดกัน โมเมนตัมคงตัว แต่พลังงานจลน์ลดลง

ถ้ามวล $m_1$ เคลื่อนที่ด้วย $u_1$ และมวล $m_2$ เคลื่อนที่ด้วย $u_2$ แล้วติดกัน ความเร็วหลังชนคือ

\[
v=\frac{m_1u_1+m_2u_2}{m_1+m_2}
\]

พลังงานจลน์ก่อนชนคือ

\[
K_i=\frac{1}{2}m_1u_1^2+\frac{1}{2}m_2u_2^2
\]

พลังงานจลน์หลังชนคือ

\[
K_f=\frac{1}{2}(m_1+m_2)v^2
\]

พลังงานที่สูญเสียคือ

\[
\Delta E_{\text{loss}}=K_i-K_f
\]

\begin{tcolorbox}[title=ข้อควรจำ]
ในการชนไม่ยืดหยุ่น พลังงานไม่ได้หายไป แต่เปลี่ยนเป็นพลังงานภายใน ความร้อน เสียง หรือการเสียรูป
\end{tcolorbox}

\subsection{ใช้โมเมนตัมร่วมกับพลังงาน}

โจทย์ชนหลายข้อแก้ด้วยสองขั้นตอน

\begin{enumerate}
    \item ช่วงการชน ใช้โมเมนตัม เพราะแรงชนเกิดในเวลาสั้น
    \item หลังชน ใช้พลังงาน ถ้าวัตถุไถลขึ้นสูงหรืออัดสปริง
\end{enumerate}

\begin{examplebox}{ตัวอย่าง: ชนแล้ววัตถุอีกก้อนแกว่งขึ้น}
มวล $m$ เคลื่อนที่ด้วยความเร็ว $u$ เข้าชนมวล $M$ ที่แขวนอยู่นิ่ง หลังชนมวล $m$ ตกลงแนวดิ่ง ทำให้โมเมนตัมแนวราบของ $m$ หลังชนเป็นศูนย์ จงหาอัตราเร็วของ $M$ หลังชนทันที

ช่วงชนสั้นมาก ใช้อนุรักษ์โมเมนตัมแนวราบ

ก่อนชน

\[
p_i=mu
\]

หลังชน มวล $m$ ไม่มีโมเมนตัมแนวราบ และมวล $M$ มีความเร็ว $V$

\[
p_f=MV
\]

ดังนั้น

\[
mu=MV
\]

\[
\boxed{V=\frac{m}{M}u}
\]

ถ้าต้องการความสูงสูงสุดหลังชน ใช้พลังงาน

\[
\frac{1}{2}MV^2=Mgh
\]

\[
h=\frac{V^2}{2g}
\]
\end{examplebox}

% =========================================================
\section{ชุดโจทย์รวมท้ายบท: คละระดับสำหรับ สอวน. ค่าย 1}

\subsection*{ระดับพื้นฐาน}

\begin{enumerate}
    \item วัตถุมวล $2\,\si{kg}$ เคลื่อนที่ด้วยความเร็ว $5\,\si{m/s}$ จงหาโมเมนตัม
    \item แรงเฉลี่ย $20\,\si{N}$ กระทำต่อวัตถุเป็นเวลา $0.5\,\si{s}$ จงหาดล
    \item วัตถุมวล $1\,\si{kg}$ เปลี่ยนความเร็วจาก $3\,\si{m/s}$ เป็น $9\,\si{m/s}$ ในทิศเดียวกัน จงหาการเปลี่ยนโมเมนตัม
    \item มวล $m$ เคลื่อนที่ด้วยความเร็ว $u$ ชนติดกับมวล $m$ ที่อยู่นิ่ง จงหาความเร็วหลังชน
\end{enumerate}

\subsection*{ระดับกลาง}

\begin{enumerate}
    \setcounter{enumi}{4}
    \item มวล $m$ เคลื่อนที่ด้วยความเร็ว $u$ ชนติดกับมวล $2m$ ที่อยู่นิ่ง จงหาพลังงานจลน์หลังชนเทียบกับก่อนชน
    \item มวล $m$ เคลื่อนที่ด้วยความเร็ว $u$ ชนยืดหยุ่นในหนึ่งมิติกับมวล $3m$ ที่อยู่นิ่ง จงหาความเร็วหลังชนของมวลทั้งสอง
    \item ลูกบอลมวล $m$ ชนกำแพงด้วยความเร็ว $u$ แล้วสะท้อนกลับด้วยความเร็ว $u$ ถ้าเวลาชนคือ $\Delta t$ จงหาขนาดแรงเฉลี่ย
    \item มวล $m$ เคลื่อนที่ด้วยความเร็ว $u$ เข้าชนสปริงที่ติดกับมวล $M$ ซึ่งอยู่นิ่งบนพื้นลื่น จงหาความเร็วของทั้งสองมวลเมื่อสปริงถูกอัดมากที่สุด
\end{enumerate}

\subsection*{ระดับยาก}

\begin{enumerate}
    \setcounter{enumi}{8}
    \item มวล $m$ เคลื่อนที่ด้วยความเร็ว $u$ ตามแกน $+x$ แล้วระเบิดเป็นสองชิ้น ชิ้นแรกมีมวล $\alpha m$ เคลื่อนที่ด้วยความเร็ว $v$ ตามแกน $+y$ จงหาเวกเตอร์ความเร็วของชิ้นที่สอง
    \item มวล $M$ เคลื่อนที่ด้วยความเร็ว $u$ เข้าชนมวล $m$ ที่อยู่นิ่งทีละลูก และติดกันไปทุกครั้ง จงหาความเร็วหลังชนลูกที่ $N$
    \item มวล $m$ เคลื่อนที่ด้วยความเร็ว $u$ ชนมวล $M$ ที่อยู่นิ่ง หลังชนมวล $m$ ไม่มีความเร็วแนวเดิม จงหาอัตราเร็วของ $M$ หลังชนทันที และความสูงสูงสุดถ้า $M$ แกว่งขึ้น
    \item ทรงกลมมวลเท่ากันสองลูกชนกันแบบยืดหยุ่น ลูกแรกมีความเร็ว $u$ และลูกที่สองอยู่นิ่ง ถ้าแนวเส้นศูนย์กลางขณะชนทำมุม $\phi$ กับความเร็วต้น จงหาขนาดความเร็วของลูกแรกหลังชน
\end{enumerate}

% =========================================================
\section{เฉลยย่อท้ายบท}

\begin{enumerate}
    \item
    \[
    p=mv=(2)(5)=10\,\si{kg.m/s}
    \]

    \item
    \[
    J=F_{\text{avg}}\Delta t=(20)(0.5)=10\,\si{N.s}
    \]

    \item
    \[
    \Delta p=m(v_f-v_i)=(1)(9-3)=6\,\si{kg.m/s}
    \]

    \item
    \[
    mu=(m+m)v
    \]
    \[
    v=\frac{u}{2}
    \]

    \item
    หลังชนติดกัน
    \[
    mu=(m+2m)v
    \]
    \[
    v=\frac{u}{3}
    \]
    พลังงานจลน์หลังชนคือ
    \[
    K_f=\frac{1}{2}(3m)\left(\frac{u}{3}\right)^2
    \]
    \[
    K_f=\frac{1}{6}mu^2
    \]
    พลังงานจลน์ก่อนชนคือ
    \[
    K_i=\frac{1}{2}mu^2
    \]
    ดังนั้น
    \[
    \frac{K_f}{K_i}=\frac{1}{3}
    \]

    \item
    สำหรับชนยืดหยุ่นและมวลที่สองอยู่นิ่ง
    \[
    v_1=\frac{m-3m}{m+3m}u=-\frac{u}{2}
    \]
    \[
    v_2=\frac{2m}{m+3m}u=\frac{u}{2}
    \]

    \item
    เลือกทิศเข้าหากำแพงเป็นบวก
    \[
    v_i=+u,\quad v_f=-u
    \]
    \[
    |\Delta p|=|m(-u)-mu|=2mu
    \]
    ดังนั้น
    \[
    F_{\text{avg}}=\frac{2mu}{\Delta t}
    \]

    \item
    เมื่อสปริงถูกอัดมากที่สุด ทั้งสองมวลมีความเร็วเท่ากัน
    \[
    mu=(m+M)V
    \]
    \[
    V=\frac{mu}{m+M}
    \]

    \item
    จากตัวอย่างในบท
    \[
    v_{2x}=\frac{u}{1-\alpha}
    \]
    \[
    v_{2y}=-\frac{\alpha v}{1-\alpha}
    \]
    ดังนั้น
    \[
    \vec{v}_2=\frac{u}{1-\alpha}\hat{i}-\frac{\alpha v}{1-\alpha}\hat{j}
    \]

    \item
    หลังชนติดกันไป $N$ ลูก มวลรวมคือ
    \[
    M+Nm
    \]
    โมเมนตัมเริ่มต้นคือ
    \[
    Mu
    \]
    ดังนั้น
    \[
    v_N=\frac{M}{M+Nm}u
    \]

    \item
    ช่วงชนใช้โมเมนตัมแนวเดิม
    \[
    mu=MV
    \]
    \[
    V=\frac{m}{M}u
    \]
    ถ้า $M$ แกว่งขึ้น
    \[
    \frac{1}{2}MV^2=Mgh
    \]
    \[
    h=\frac{V^2}{2g}=\frac{m^2u^2}{2M^2g}
    \]

    \item
    แตกความเร็วของลูกแรกเป็นแนวเส้นศูนย์กลางและแนวสัมผัส
    \[
    u_n=u\cos\phi
    \]
    \[
    u_t=u\sin\phi
    \]
    สำหรับมวลเท่ากัน ชนยืดหยุ่น และลูกที่สองอยู่นิ่ง องค์ประกอบแนวเส้นศูนย์กลางของลูกแรกถ่ายไปให้ลูกที่สอง ส่วนลูกแรกเหลือองค์ประกอบแนวสัมผัส
    \[
    \boxed{v_1'=u\sin\phi}
    \]
\end{enumerate}

% =========================================================
\section{สรุปสูตรสำคัญ}

\begin{tcolorbox}[title=สรุปบท: โมเมนตัมเชิงเส้นและการชน]
โมเมนตัม

\[
\vec{p}=m\vec{v}
\]

กฎข้อที่สองในรูปโมเมนตัม

\[
\sum \vec{F}=\frac{d\vec{p}}{dt}
\]

ดล

\[
\vec{J}=\int \vec{F}\,dt
\]

\[
\vec{J}=\vec{F}_{\text{avg}}\Delta t
\]

ทฤษฎีบทดลและโมเมนตัม

\[
\vec{J}_{\text{net}}=\Delta \vec{p}
\]

อนุรักษ์โมเมนตัม

\[
\vec{P}_i=\vec{P}_f
\]

\[
\sum_i m_i\vec{v}_{i,\text{before}}
=
\sum_i m_i\vec{v}_{i,\text{after}}
\]
\end{tcolorbox}

\begin{tcolorbox}[title=สรุปบท: โมเมนตัมเชิงเส้นและการชน ต่อ]
ชนติดกัน

\[
v=\frac{m_1u_1+m_2u_2}{m_1+m_2}
\]

ชนยืดหยุ่นหนึ่งมิติ

\[
v_1=
\frac{m_1-m_2}{m_1+m_2}u_1
+
\frac{2m_2}{m_1+m_2}u_2
\]

\[
v_2=
\frac{2m_1}{m_1+m_2}u_1
+
\frac{m_2-m_1}{m_1+m_2}u_2
\]

ศูนย์กลางมวล

\[
\vec{r}_{\text{cm}}=\frac{\sum_i m_i\vec{r}_i}{M_{\text{tot}}}
\]

\[
\vec{P}_{\text{tot}}=M_{\text{tot}}\vec{v}_{\text{cm}}
\]

การชนสองมิติ

\[
\sum p_{xi}=\sum p_{xf}
\]

\[
\sum p_{yi}=\sum p_{yf}
\]
\end{tcolorbox}

% =========================================================
\section{ข้อสอบจริง สอวน. ที่ใช้ความรู้บทโมเมนตัมและการชน}

\begin{tcolorbox}[title=คำแนะนำการใส่ภาพข้อสอบ]
ให้ตัดภาพข้อสอบจริงแล้วบันทึกตามชื่อไฟล์ที่กำหนดไว้ในแต่ละข้อ ถ้ายังไม่มีไฟล์ ภาพจะไม่แสดง แต่เอกสารยังคอมไพล์ได้
\end{tcolorbox}

% =========================================================
\subsection{ปี 2560 ตอนที่ 1 ข้อ 5: ลูกบอลกระดอนและแรงเฉลี่ยจากพื้น}

\vspace{1em}

\IfFileExists{img/posn60p1q5.png}{
\begin{center}
    \includegraphics[width=1\linewidth]{img/posn60p1q5.png}
\end{center}
}{
\begin{tcolorbox}[title=ตำแหน่งภาพที่แนะนำ]
ให้ตัดภาพข้อสอบปี 2560 ตอนที่ 1 ข้อ 5 แล้วบันทึกเป็น

\[
\texttt{img/posn60p1q5.png}
\]
\end{tcolorbox}
}

\begin{examplebox}{เฉลยละเอียด}
ลูกบอลทั้งสองมีมวลเท่ากัน ปล่อยจากความสูงเท่ากัน ดังนั้นก่อนกระทบพื้น ลูกบอลทั้งสองมีอัตราเร็วเท่ากัน ให้เป็น

\[
u
\]

หลังชน ลูกที่ 1 กระดอนสูงกว่าลูกที่ 2 ดังนั้นอัตราเร็วหลังชนของลูกที่ 1 มากกว่า

\[
v_1>v_2
\]

เลือกทิศขึ้นเป็นบวก ก่อนชนความเร็วคือ

\[
-u
\]

หลังชนความเร็วคือ

\[
+v
\]

การเปลี่ยนโมเมนตัมของลูกบอลคือ

\[
\Delta p=m(v-(-u))
\]

\[
\Delta p=m(v+u)
\]

ในช่วงเวลาสัมผัสพื้นเท่ากัน แรงเฉลี่ยจากพื้นแปรตามดลที่พื้นให้ลูกบอล

เพราะลูกที่ 1 มี $v$ หลังชนมากกว่า จึงมีการเปลี่ยนโมเมนตัมมากกว่า

ดังนั้นแรงเฉลี่ยจากพื้นต่อลูกที่ 1 มากกว่าแรงเฉลี่ยจากพื้นต่อลูกที่ 2

\[
\boxed{F_1>F_2}
\]
\end{examplebox}

% =========================================================
\subsection{ปี 2560 ตอนที่ 1 ข้อ 6: ชนแล้วลูกตุ้มแกว่งขึ้น}

\vspace{1em}

\IfFileExists{img/posn60p1q6.png}{
\begin{center}
    \includegraphics[width=1\linewidth]{img/posn60p1q6.png}
\end{center}
}{
\begin{tcolorbox}[title=ตำแหน่งภาพที่แนะนำ]
ให้ตัดภาพข้อสอบปี 2560 ตอนที่ 1 ข้อ 6 แล้วบันทึกเป็น

\[
\texttt{img/posn60p1q6.png}
\]
\end{tcolorbox}
}

\begin{examplebox}{เฉลยละเอียด}
มวลที่พุ่งเข้าชนคือ

\[
m=0.20\,\si{kg}
\]

มีความเร็วต้น

\[
u=3.0\,\si{m/s}
\]

ทรงกลมที่แขวนมีมวล

\[
M=1.0\,\si{kg}
\]

หลังชน มวล $m$ ตกลงตรง ๆ ในแนวดิ่ง ดังนั้นโมเมนตัมแนวราบของ $m$ หลังชนเป็นศูนย์

ช่วงการชนใช้อนุรักษ์โมเมนตัมแนวราบ

\[
mu=MV
\]

ดังนั้น

\[
V=\frac{m}{M}u
\]

แทนค่า

\[
V=\frac{0.20}{1.0}(3.0)
\]

\[
V=0.60\,\si{m/s}
\]

หลังชน ทรงกลมมวล $M$ แกว่งขึ้น โดยพลังงานจลน์เปลี่ยนเป็นพลังงานศักย์โน้มถ่วง

\[
\frac{1}{2}MV^2=Mgh
\]

ตัด $M$

\[
h=\frac{V^2}{2g}
\]

แทนค่า $g=9.8\,\si{m/s^2}$

\[
h=\frac{(0.60)^2}{2(9.8)}
\]

\[
h=\frac{0.36}{19.6}
\]

\[
h\approx 0.018\,\si{m}
\]

ดังนั้น

\[
\boxed{h\approx 0.018\,\si{m}}
\]
\end{examplebox}

% =========================================================
\subsection{ปี 2561 ข้อ 6: นิวตรอนชนดิวเทอรอนแบบยืดหยุ่น}

\vspace{1em}

\IfFileExists{img/posn61q6.png}{
\begin{center}
    \includegraphics[width=1\linewidth]{img/posn61q6.png}
\end{center}
}{
\begin{tcolorbox}[title=ตำแหน่งภาพที่แนะนำ]
ให้ตัดภาพข้อสอบปี 2561 ข้อ 6 แล้วบันทึกเป็น

\[
\texttt{img/posn61q6.png}
\]
\end{tcolorbox}
}

\begin{examplebox}{เฉลยละเอียด}
ให้นิวตรอนมีมวล

\[
m_1=m
\]

และดิวเทอรอนมีมวล

\[
m_2=2m
\]

ดิวเทอรอนอยู่นิ่งก่อนชน

\[
u_2=0
\]

นิวตรอนมีความเร็วต้น

\[
u_1=u
\]

สำหรับการชนยืดหยุ่นหนึ่งมิติ

\[
v_1=\frac{m_1-m_2}{m_1+m_2}u_1
\]

แทนค่า

\[
v_1=\frac{m-2m}{m+2m}u
\]

\[
v_1=-\frac{u}{3}
\]

ดังนั้นพลังงานจลน์ของนิวตรอนหลังชนเทียบกับก่อนชนคือ

\[
\frac{K_f}{K_i}
=
\frac{\frac{1}{2}m\left(\frac{u}{3}\right)^2}
{\frac{1}{2}mu^2}
\]

\[
\boxed{\frac{K_f}{K_i}=\frac{1}{9}}
\]
\end{examplebox}

% =========================================================
\subsection{ปี 2566 ข้อ 3: การระเบิดและทิศของชิ้นส่วน}

\vspace{1em}

\IfFileExists{img/posn66q3.png}{
\begin{center}
    \includegraphics[width=1\linewidth]{img/posn66q3.png}
\end{center}
}{
\begin{tcolorbox}[title=ตำแหน่งภาพที่แนะนำ]
ให้ตัดภาพข้อสอบปี 2566 ข้อ 3 แล้วบันทึกเป็น

\[
\texttt{img/posn66q3.png}
\]
\end{tcolorbox}
}

\begin{examplebox}{เฉลยละเอียด}
มวล $m$ เคลื่อนที่ด้วยความเร็ว $u$ ตามแกน $+x$ ก่อนระเบิด

\[
\vec{p}_i=mu\hat{i}
\]

หลังระเบิด ชิ้นแรกมีมวล

\[
\alpha m
\]

และเคลื่อนที่ด้วยความเร็ว $v$ ตามแกน $+y$

\[
\vec{p}_1=\alpha mv\hat{j}
\]

ชิ้นที่สองมีมวล

\[
(1-\alpha)m
\]

ให้ความเร็วของชิ้นที่สองเป็น

\[
\vec{v}_2=v_{2x}\hat{i}+v_{2y}\hat{j}
\]

ใช้อนุรักษ์โมเมนตัม

\[
mu\hat{i}
=
\alpha mv\hat{j}
+
(1-\alpha)m(v_{2x}\hat{i}+v_{2y}\hat{j})
\]

แยกแกน $x$

\[
mu=(1-\alpha)mv_{2x}
\]

\[
v_{2x}=\frac{u}{1-\alpha}
\]

แยกแกน $y$

\[
0=\alpha mv+(1-\alpha)mv_{2y}
\]

\[
v_{2y}=-\frac{\alpha v}{1-\alpha}
\]

ดังนั้นมุมที่ชิ้นที่สองทำกับแกน $+x$ ลงด้านล่างคือ

\[
\tan\theta=\frac{|v_{2y}|}{v_{2x}}
\]

\[
\boxed{\tan\theta=\frac{\alpha v}{u}}
\]
\end{examplebox}

% =========================================================
\subsection{ปี 2566 ข้อ 4: ก้อนน้ำชนกำแพงและแรงเฉลี่ย}

\vspace{1em}

\IfFileExists{img/posn66q4.png}{
\begin{center}
    \includegraphics[width=1\linewidth]{img/posn66q4.png}
\end{center}
}{
\begin{tcolorbox}[title=ตำแหน่งภาพที่แนะนำ]
ให้ตัดภาพข้อสอบปี 2566 ข้อ 4 แล้วบันทึกเป็น

\[
\texttt{img/posn66q4.png}
\]
\end{tcolorbox}
}

\begin{examplebox}{เฉลยละเอียด}
ก้อนน้ำเหลวทรงกลมมีมวล

\[
m
\]

รัศมี

\[
r
\]

และเคลื่อนที่เข้ากำแพงด้วยความเร็ว

\[
v
\]

เมื่อก้อนน้ำชนกำแพง โมเมนตัมแนวเข้ากำแพงถูกลดจาก

\[
mv
\]

เป็นศูนย์โดยประมาณ ดังนั้นขนาดการเปลี่ยนโมเมนตัมคือ

\[
\Delta p=mv
\]

เวลาที่ก้อนน้ำทั้งก้อนใช้ในการชนกำแพงประมาณเท่ากับเวลาที่ความยาวเส้นผ่านศูนย์กลาง $2r$ เคลื่อนผ่านกำแพง

\[
\Delta t=\frac{2r}{v}
\]

แรงเฉลี่ยจากกำแพงคือ

\[
F_{\text{avg}}=\frac{\Delta p}{\Delta t}
\]

แทนค่า

\[
F_{\text{avg}}
=
\frac{mv}{\frac{2r}{v}}
\]

\[
\boxed{
F_{\text{avg}}=\frac{1}{2}\frac{mv^2}{r}
}
\]
\end{examplebox}

% =========================================================
\subsection{ปี 2566 ข้อ 6: ชนติดกันทีละลูก}

\vspace{1em}

\IfFileExists{img/posn66q6.png}{
\begin{center}
    \includegraphics[width=1\linewidth]{img/posn66q6.png}
\end{center}
}{
\begin{tcolorbox}[title=ตำแหน่งภาพที่แนะนำ]
ให้ตัดภาพข้อสอบปี 2566 ข้อ 6 แล้วบันทึกเป็น

\[
\texttt{img/posn66q6.png}
\]
\end{tcolorbox}
}

\begin{examplebox}{เฉลยละเอียด}
มวล $M$ มีความเร็วต้น

\[
u
\]

และเคลื่อนที่เข้าชนมวลเล็ก $m$ ทีละลูก โดยชนแล้วติดกันไปทุกครั้ง

ก่อนชนทั้งหมด โมเมนตัมรวมคือ

\[
P_i=Mu
\]

หลังชนลูกที่ $N$ แล้ว มวลรวมที่เคลื่อนที่ติดกันคือ

\[
M+Nm
\]

ให้ความเร็วหลังชนลูกที่ $N$ เป็น

\[
v_N
\]

ใช้อนุรักษ์โมเมนตัม

\[
Mu=(M+Nm)v_N
\]

ดังนั้น

\[
\boxed{
v_N=\frac{M}{M+Nm}u
}
\]
\end{examplebox}

% =========================================================
\subsection{ปี 2568 ข้อ 2: ทรงกลมมวลเท่ากันชนเฉียงแบบยืดหยุ่น}

\vspace{1em}

\IfFileExists{img/posn68q2.png}{
\begin{center}
    \includegraphics[width=1\linewidth]{img/posn68q2.png}
\end{center}
}{
\begin{tcolorbox}[title=ตำแหน่งภาพที่แนะนำ]
ให้ตัดภาพข้อสอบปี 2568 ข้อ 2 แล้วบันทึกเป็น

\[
\texttt{img/posn68q2.png}
\]
\end{tcolorbox}
}

\begin{examplebox}{เฉลยละเอียด}
ทรงกลม $A$ และ $B$ มีมวลเท่ากัน และผิวเกลี้ยงลื่น

ทรงกลม $A$ เคลื่อนที่ด้วยความเร็วต้น

\[
u
\]

ทรงกลม $B$ อยู่นิ่ง และชนกันแบบยืดหยุ่น

จากรูป แนวเส้นศูนย์กลางขณะชนทำมุม $30^\circ$ กับแนวความเร็วต้นของ $A$

แตกความเร็วของ $A$ เป็นสองแนว

แนวเส้นศูนย์กลาง

\[
u_n=u\cos30^\circ
\]

แนวสัมผัส

\[
u_t=u\sin30^\circ
\]

เพราะผิวลื่น แรงดลระหว่างชนมีเฉพาะตามแนวเส้นศูนย์กลาง ดังนั้นองค์ประกอบแนวสัมผัสของ $A$ ไม่เปลี่ยน

สำหรับมวลเท่ากัน ชนยืดหยุ่น และ $B$ อยู่นิ่ง องค์ประกอบแนวเส้นศูนย์กลางของ $A$ จะถ่ายไปให้ $B$

ดังนั้นหลังชน $A$ เหลือเฉพาะองค์ประกอบแนวสัมผัส

\[
v_A'=u_t
\]

\[
v_A'=u\sin30^\circ
\]

\[
\boxed{v_A'=\frac{u}{2}}
\]
\end{examplebox}

% =========================================================
\subsection{ปี 2568 ข้อ 3: มวลชนสปริงที่ติดกับมวลอีกก้อน}

\vspace{1em}

\IfFileExists{img/posn68q3.png}{
\begin{center}
    \includegraphics[width=1\linewidth]{img/posn68q3.png}
\end{center}
}{
\begin{tcolorbox}[title=ตำแหน่งภาพที่แนะนำ]
ให้ตัดภาพข้อสอบปี 2568 ข้อ 3 แล้วบันทึกเป็น

\[
\texttt{img/posn68q3.png}
\]
\end{tcolorbox}
}

\begin{examplebox}{เฉลยละเอียด}
มวล $m$ เคลื่อนที่ด้วยความเร็ว

\[
u
\]

เข้าชนและอัดสปริงที่ติดกับมวล $M$ ซึ่งอยู่นิ่งก่อนชน

ช่วงที่สปริงถูกอัดมากที่สุด มวล $m$ กับมวล $M$ เข้าใกล้กันไม่ได้อีกชั่วขณะ ดังนั้นทั้งสองมีความเร็วเท่ากัน ให้เป็น

\[
V
\]

ไม่มีแรงภายนอกในแนวราบ จึงใช้อนุรักษ์โมเมนตัม

ก่อนชน

\[
P_i=mu
\]

ขณะสปริงถูกอัดมากที่สุด

\[
P_f=(m+M)V
\]

ดังนั้น

\[
mu=(m+M)V
\]

\[
\boxed{
V=\frac{mu}{m+M}
}
\]

นี่คือความเร็วของ $m$ ที่จังหวะที่ $m$ เข้าใกล้ $M$ มากที่สุด
\end{examplebox}